\chapter{Introducción a la teoría espectral}
\label{ch:teoria-espectral}
\index{Valor propio}\index{Función propia}\index{Operador compacto}\index{Operador autoadjunto}\index{Espectro}\index{Descomposición espectral}

\section{Valores propios}

La teoría espectral busca generalizar, a operadores en dimensiones infinitas, la descomposición diagonal de matrices. David Hilbert, en sus trabajos sobre ecuaciones integrales (1904--1910), descubrió que muchos operadores importantes admiten una ``descomposición espectral'' análoga a la diagonalización, aunque los ``eigenvalores'' podían formar un espectro continuo. Esta teoría constituye uno de los puentes más profundos entre el análisis funcional y las aplicaciones físicas.

\begin{definition}
\label{def:espectro}
Sea $E$ un espacio de Banach y $T\colon E\to E$ un operador lineal acotado. El \emph{espectro} $\sigma(T)$ es el conjunto de $\lambda\in\mathbb{C}$ tales que $T-\lambda I$ no es invertible con inversa acotada. Los $\lambda$ para los cuales $T-\lambda I$ no es inyectivo son los \emph{valores propios} (o \emph{eigenvalores}).
\end{definition}

\begin{example}
\label{ej:multiplicacion}
Sea $T\colon L^2([0,1])\to L^2([0,1])$ definido por $Tf(x)=xf(x)$. Entonces $T$ no tiene valores propios, pero su espectro es $[0,1]$ (espectro continuo).
\end{example}

\section{Funciones propias}

\begin{definition}
\label{def:funcion-propia}
Si $\lambda$ es un valor propio de $T$, un vector no nulo $v$ tal que $Tv=\lambda v$ se denomina \emph{función propia} (o \emph{eigenvector}).
\end{definition}

\begin{example}
\label{ej:derivada}
El operador derivada $D\colon C^1([0,2\pi])\to C([0,2\pi])$ tiene valores propios $\lambda_n=in$ con funciones propias $e^{inx}$.
\end{example}

\section{Operadores compactos}

\begin{theorem}[Riesz--Schauder]
\label{teo:riesz-schauder}
Si $T$ es un operador compacto en un espacio de Banach, su espectro es a lo sumo numerable, $0$ es el único posible punto de acumulación, y todo $\lambda\neq 0$ en el espectro es un valor propio de multiplicidad finita.
\end{theorem}

\section{Operadores autoadjuntos (idea)}

\begin{definition}
\label{def:autoadjunto}
En un espacio de Hilbert, un operador $T$ es \emph{autoadjunto} si $\langle Tx,y\rangle=\langle x,Ty\rangle$ para todo $x,y$.
\end{definition}

\begin{theorem}[Espectral para autoadjuntos compactos]
\label{teo:espectral-compacto}
Si $T$ es autoadjunto y compacto, existe una base ortonormal de $H$ formada por funciones propias de $T$. Los valores propios son reales y tienden a $0$.
\end{theorem}

\section{Descomposición espectral (motivación)}

El teorema espectral general afirma que todo operador autoadjunto (no necesariamente compacto) admite una ``medida espectral'' $E$ tal que
\[
  T=\int_{\sigma(T)}\lambda\,dE(\lambda).
\]
Esta fórmula generaliza la diagonalización $A=PDP^{-1}$ de matrices simétricas.

\section{Aplicaciones a ecuaciones diferenciales}

\begin{example}[Problema de Sturm--Liouville]
\label{ej:sturm-liouville}
La ecuación $-(py')'+qy=\lambda w y$ en $[a,b]$ con condiciones de contorno define un operador autoadjunto cuyos valores propios forman una sucesión creciente $\lambda_n\to\infty$, y las funciones propias constituyen una base ortogonal de $L^2_w([a,b])$.
\end{example}

\begin{exercise}
Demuestre que los valores propios de un operador autoadjunto son reales.
\end{exercise}

\begin{exercise}
Sea $T\colon\ell^2\to\ell^2$ el operador shift derecha $T(x_1,x_2,\dots)=(0,x_1,x_2,\dots)$. Determine su espectro.
\end{exercise}

\begin{exercise}
Demuestre que si $T$ es compacto y autoadjunto, entonces $\|T\|$ es igual al máximo de los valores absolutos de sus valores propios.
\end{exercise}

\begin{problem}
Formule y demuestre un teorema de alternativa de Fredholm para operadores compactos: la ecuación $u-Tu=f$ tiene solución si y solo si $f$ es ortogonal al núcleo de $I-T^*$.
\end{problem}
